% ==============================================================================
% CEBE Interdisciplinary Fellowship -- CV template
%
% Max 2 pages per person. Copy this file once per applicant, e.g.:
%   cv-pi.tex, cv-cosupervisor1.tex, cv-cosupervisor2.tex, cv-candidate.tex
% and compile each separately, then merge with main.tex into the single
% submission PDF (see README.md).
%
% Required minimum content per the call:
%   - Name, title and affiliation
%   - PhD year
%   - Short statement on previous research achievements and relevant
%     impacts (e.g. policy, industrial, or public arenas)
%   - Number (not a name list) of current and former PhD students,
%     postdocs, and master's students
%   - 10 publications: 5 most important (last 15 years) + 5 recent
%     project-relevant ones (last 5 years)
%   - Link to Scopus profile or ORCID
% ==============================================================================

% ==============================================================================
% Instantiated from templates/cv-template.tex on 2026-09-02 at Giuseppe's
% request. Name/title/affiliation below are pre-filled from AGENTS.md's
% "Team" section where known; everything else (PhD year, achievements,
% supervision counts, publications, ORCID/Scopus) is still a placeholder.
% ==============================================================================

\documentclass[12pt,a4paper]{article}

\usepackage[a4paper,margin=2.5cm]{geometry}
\usepackage[T1]{fontenc}
\usepackage[utf8]{inputenc}
\usepackage{mathptmx}
\usepackage{setspace}
\singlespacing
\usepackage{enumitem}
\usepackage{xcolor}
\usepackage[colorlinks=true,linkcolor=blue,urlcolor=blue]{hyperref}

\newcommand{\pagemark}[1]{\typeout{PAGEMARK: #1 -- page \thepage}}

% Expected-length annotation (matches proposal/main-template.tex). Set
% \templatenotesfalse to hide before the final submission build.
\newif\iftemplatenotes
\templatenotestrue
\newcommand{\lengthnote}[1]{%
  \iftemplatenotes
    {\small\itshape\color{gray!70!black} [Target length: #1]}\par\vspace{4pt}
  \fi
}

\pagenumbering{gobble} % CVs typically don't need page numbers in a merged submission; remove if you want them

\begin{document}
\pagemark{CV Giuseppe Abbiati (PI) -- start}
\lengthnote{max 2 pages per person (per the call)}

\begin{center}
% Full name inferred from git/GitHub identity (giuseppe.abbiati@gmail.com,
% github.com/giuseppeabbiati1984) -- confirm/correct if needed.
{\Large\bfseries Giuseppe Abbiati}\par\vspace{2pt}
[Associate Professor in Structural Mechanics]\par
Dept. of Civil and Architectural Engineering (CAE), Aarhus University, Denmark
\end{center}

\vspace{0.3cm}
\noindent\textbf{ORCID / Scopus:} \url{[link]}

\vspace{0.5cm}
\noindent\textbf{Education}
PhD
MSc
BSc

\vspace{0.5cm}
\noindent\textbf{Employment}

\vspace{0.5cm}
\noindent\textbf{Research achievements and impact}

[Short statement on previous research achievements and relevant impacts,
including e.g. policy, industrial, or public-arena impact. A few
sentences, not a full narrative CV.]

\vspace{0.5cm}
\noindent\textbf{Supervision record}

Current: 5 PhD students, 2 postdocs, 7master's students.\\
Total: 7 PhD students, 4 postdocs, 45 master's students.

\vspace{0.5cm}
\noindent\textbf{Selected publications}

\noindent\textit{Five most important (last 15 years):}
\begin{enumerate}%[leftmargin=1.2cm]
  \item Scussolini, L., Abbiati, G., \& Ceravolo, R. (2026). Identification of Volterra time-varying parameters in dynamical systems under random excitation. Journal of Sound and Vibration, 645, 120128, \url{10.1016/j.jsv.2026.120128}
  \item Pizarro, D., M. Kovarbašić, G. Abbiati, and B. Stojadinović, (2025). “Multi‐Axial Subassemblage Testing System for Hybrid Simulation with Large‐Scale Structural Components.” Earthquake Engineering \& Structural Dynamics, 54(8):2084-2105,\url{10.1002/eqe.4351}.
  \item Stamenov, D., Abbiati, G., \& Sauder, T. (2025). Numerical estimation of generalized frequency response functions from time series data using NARX. Mechanical Systems and Signal Processing, 239, 113278, \url{10.1016/j.ymssp.2025.113278}. 
  \item G. Abbiati, M. Broccardo, I. Abdallah, S. Marelli, F. Paolacci, 2021. Seismic Fragility Analysis based on Artificial Ground Motions and Surrogate Modeling of Validated Structural Simulators, Earthquake Engineering and Structural Dynamics, 50(9): 2314-2333, ISSN 00988847, \url{10.1002/eqe.3448}.
  \item Abbiati G., R. Ceravolo, C. Surace (2015). Time-dependent Estimators for On-Line Monitoring of Full-Scale Structures under Ambient Excitation, Mechanical Systems and Signal Processing, 60:166-181, ISSN 08883270, \url{10.1016/j.ymssp.2014.10.018}.
\end{enumerate}

\noindent\textit{Five recent, project-relevant (last 5 years):}
\begin{enumerate}%[leftmargin=1.2cm]
  \item P. Talasila, D. Tcherniak, A. M.D. Jensen, S Mahato, J. V. Medom, M. D. Ulriksen, G. Abbiati, A. Schörghofer-Queiroz, P. G. Larsen, L. Damkilde, (2026). A Digital Twin Platform for Structural Health Monitoring, Computer-Aided Civil and Infrastructure Engineering, 50:100086, ISSN 1093-9687, DOI 10.1016/j.cacaie.2026.100086.
  \item Hansen, S. T., Thule, C., Gomes, C., Lausdahl, K. G., Madsen, F. P., Abbiati, G., & Larsen, P. G. (2024). Co-simulation at different levels of expertise with Maestro2. Journal of Systems and Software, 209, 111905. https://doi.org/10.1016/j.jss.2023.111905
  \item Abbiati, G., C. Gomes, M. Sandberg, Z. Kazemi, S. T. Hansen, and P. G. Larsen. “Modelling for Digital Twins.” In The Engineering of Digital Twins, edited by J. Fitzgerald, C. Gomes, and P. G. Larsen, 89–127. Cham: Springer International Publishing, 2024. https://doi.org/10.1007/978-3-031-66719-0_5.
  \item Oakes, B, C Gomes, G Abbiati, E Kamburjan, E E. Baş, and S Engelsgaard. “Towards Ontological Service-Driven Engineering of Digital Twins.” In 1st International Conference on Engineering Digital Twins (EDTconf 2024). Linz, Austria: to appear, 2024.
  \item Baş, E. E., Abbiati, G., Gonçalves Gomes, C. Â., Jassmann, U., & Larsen, P. G. (2024). Digitalization of Large-Scale Testing Facilities for the Wind Industry: DIGIT-BENCH Digital Twin. Journal of Physics: Conference Series, 2767(4), 042033. https://doi.org/10.1088/1742-6596/2767/4/042033

\end{enumerate}

\pagemark{CV Giuseppe Abbiati (PI) -- end}
\end{document}
